\documentclass[11pt]{article}
\usepackage{amsmath,amssymb,amsthm}
\newtheorem{theorem}{Theorem}
\newtheorem{proposition}{Proposition}
\title{The Exact Affine Span of Butterfly Permutations}
\author{}
\date{November 23, 2026}
\begin{document}
\maketitle

\paragraph{Status.}
The affine span and the resulting independent diagonal readouts below have
complete written proofs and exhaustive exact checks through eight wires.
They do not give an admissible head decoder or an asymptotic complexity
bound. None of this note is Lean formalized. The final obstruction concerns
exact single-kernel readouts with fixed residues only.

Put $H=2^d$, and label the wires by $d$ bits numbered from $0$ to $d-1$.
At layer $k$, the network independently switches the two wires differing
in bit $k$, for every fixed lower tuple $s$ and higher tuple $g$.
Write $\mathcal P_d$ for the resulting permutation matrices, with
$P_\pi(h,y)=\mathbf1_{\{\pi(h)=y\}}$.
There are $Hd/2$ switch bits. Their independent rectangle readouts and the
injectivity of the network parametrization are proved in
\texttt{butterfly-rectangle-decoder.tex}.

Define input and output Haar vectors by
\[
 u_{k,g}(h)=\mathbf1_{\{h>>(k+1)=g\}}(-1)^{h_k},
 \qquad
 v_{k,s}(y)=\mathbf1_{\{y\bmod 2^k=s\}}(-1)^{y_k}.
\]
Let $U_k$ and $V_k$ be their spans as $g$ and $s$ vary.
Together with the constant vector, the $U_k$ form an orthogonal
 decomposition of $\mathbb R^H$, and the $V_k$ do likewise in the reverse
bit order. Their dimensions are
\[
 \dim U_k=H/2^{k+1},\qquad \dim V_k=2^k.
\]
All inner products are unnormalized coordinate sums.

\begin{theorem}[Exact affine difference space]
For any $P_0\in\mathcal P_d$, the linear span of
$\{P-P_0:P\in\mathcal P_d\}$ is
\[
 \mathcal S_d=\bigoplus_{0\leq a\leq b<d}U_a\otimes V_b.
\]
In particular its dimension is $H^2-H-Hd/2$.
\end{theorem}
\begin{proof}
Every permutation matrix has row and column sums one, so difference
matrices have zero constant tensor components. Consider $a>b$.
For an input block of size $2^{a+1}$ defining $u_{a,g}$, each of its
halves has fixed bits at positions at least $a$. The first $a$ layers
biject its lower $a$ bits, and all remaining layers preserve those
output bits. Each half therefore contains exactly one image with each
prescribed lower $a$ tuple. Since $v_{b,s}$ depends only on these lower
bits, its sum over either half is the same, in fact zero. Hence
$\langle u_{a,g}\otimes v_{b,s},P\rangle=0$ for every $P$.
This proves the inclusion in the asserted upper triangular tensor sum.

For the reverse inclusion, change just one layer-$k$ switch, leaving
all other switches fixed. Its two incoming original wires $h,h'$ have
highest differing bit $k$, and its two final outputs $y,y'$ have lowest
differing bit $k$. The difference of the two permutation matrices is,
up to sign,
\[
 (e_h-e_{h'})\otimes(e_y-e_{y'}).
\]
Every pair of this form is realizable. Earlier layers can independently
route the two chosen inputs to any common lower tuple, because they
lie in distinct fixed-bit-$k$ halves. To see the routing assertion,
follow the input and set its switch at each earlier layer to assign
the desired next bit. Later layers can likewise independently assign
the two desired higher output tuples, because the two wires have
different bit $k$ and later switches preserve it.

Differences between all pairs with highest differing bit $k$ span the
space of vectors with zero sum on each dyadic block of size $2^{k+1}$.
Indeed, within each such block the cross-half pairs form a connected
complete bipartite graph; its edge differences span all zero-sum
vectors on that block. This space is $U_0\oplus\cdots\oplus U_k$.
The analogous output differences span $V_k\oplus\cdots\oplus V_{d-1}$.
Since the input and output pairs above can be selected independently,
$\mathcal S_d$ contains
\[
 (U_0\oplus\cdots\oplus U_k)\otimes
 (V_k\oplus\cdots\oplus V_{d-1})
\]
for every $k$. These spaces together contain every tensor block $a\leq b$.
Finally,
\[
 \sum_{a=0}^{d-1}\frac{H}{2^{a+1}}
       \sum_{b=a}^{d-1}2^b
 =\sum_{a=0}^{d-1}\frac{H}{2^{a+1}}(H-2^a)
 =H^2-H-Hd/2.
\]
\end{proof}

\begin{proposition}[Independent diagonal Haar signs]
If $\theta_{k,s,g}\in\{0,1\}$ is the setting of a switch, then
\[
 \langle u_{k,g}\otimes v_{k,s},P_\pi\rangle
       =2-4\theta_{k,s,g}.
\]
Thus all $Hd/2$ diagonal Haar tensor signs are independently realized.
\end{proposition}
\begin{proof}
Let $A$ be the bit-$k$ zero half of input block $g$, and let $B$ be
the bit-$k$ one half of output residue class $s$. Write
\[
 \mathbf1_A=(\mathbf1_{\text{input block }g}+u_{k,g})/2,
 \qquad
 \mathbf1_B=(\mathbf1_{\text{output residue }s}-v_{k,s})/2.
\]
The input block indicator has only a constant component and input
levels strictly above $k$; the output residue indicator has only a
constant component and output levels strictly below $k$.
Consequently the projection of
$\mathbf1_A\otimes\mathbf1_B$ onto $\mathcal S_d$ is
$-u_{k,g}\otimes v_{k,s}/4$.
The rectangle identity gives
$\langle\mathbf1_A\otimes\mathbf1_B,P_\pi\rangle=\theta_{k,s,g}$.
For the all-zero switch configuration the permutation is the identity,
the rectangle count is zero, and the Haar tensor inner product is two:
its two nonzero diagonal contributions both equal one. Subtracting this
base configuration proves the formula.
\end{proof}

\paragraph{What quotienting by the invariants permits.}
Let $a,b\in\mathbb R^H$ and consider the assignment cost
$\sum_h a_h b_{\pi(h)}$.
Its nonconstant restriction to butterfly permutations is completely
specified by the tensor blocks
\[
 \operatorname{proj}_{U_i}(a)\otimes
 \operatorname{proj}_{V_j}(b),\qquad i\leq j.
\]
Therefore an exact scalar multiple of one diagonal Haar readout, up to
an assignment-independent constant, requires these blocks to vanish
except at the prescribed pair $(k,k)$; at that pair they must be the
prescribed nonzero multiple of $u_{k,g}\otimes v_{k,s}$.
This is an equality criterion, not a necessary criterion for merely
matching signs on the finite permutation family.

\begin{proposition}[A restricted exterior-kernel obstruction]
Fix distinct real poles $\xi_h$, take $t>\max_h\xi_h$, and put
$a_h=(t-\xi_h)^{-p}$ for any positive integer $p$.
If $H\geq4$, no choice of $b$, nonzero scalar, and additive constant
makes the cost $\sum_h a_h b_{\pi(h)}$ exactly equal to any single
switch readout on all butterfly permutations.
\end{proposition}
\begin{proof}
The projection of $a$ onto $U_0$ has a nonzero coefficient on every
adjacent-pair Haar vector: the two distinct poles give unequal kernel
values. To obtain a diagonal level-$k$ readout for $k>0$, the output
level-$k$ component of $b$ must be nonzero; its tensor with the nonzero
input level-zero component then gives a forbidden $(0,k)$ block.
For $k=0$, the output level-zero space has dimension one, so the input
level-zero component would have to be supported on one input pair.
There are at least two pairs, and all its pair coefficients are nonzero.
\end{proof}

This obstruction does not apply to target-dependent residues, sums of
several exterior kernels, approximate or sign-only readouts, or the full
admissible head model. Hierarchically separated pole digits can suppress
some lower input levels, but do not by themselves isolate one location
within the desired level. Producing simultaneous strict sign readouts
with a superconstant row-to-head ratio remains open.

\paragraph{Verification.}
The companion exact checker exhausts the networks on $H=2,4,8$ wires.
It checks every diagonal identity and every vanishing lower tensor
block using integers. The affine span ranks are $1,8,44$, certified
by row reduction modulo the prime $1000003$; the written upper bound
then makes these modular lower bounds exact over the rationals.
No numerical feasibility result is promoted to a universal theorem.
\end{document}
